\subsection{Directional disjoint group-peak proximal update}\label{sec:directional}
The second proximal path prices the largest weighted departure within each
disjoint group, using \eqref{eq:group-peak}. It permits unequal upper and lower
scales and keeps the original interval-row representation in its principal
system. Existing directional row splitting and compact epigraphs express the
same objective, so row preservation is a representation property rather than
an expressibility or priority result. Coordinatewise orthant monotonicity and
weighted threshold calculations have established antecedents~\cite{orthant,weighted}.

For a box $C$ and residual penalty $\Psi$, assume $\Psi$ is proper, closed,
convex, nonnegative, minimized at zero, and nonincreasing under coordinate
contractions $D=\diag(d_i)$ with $0\leq d_i\leq1$. For fixed diagonal $H\succ0$,
the exact identity is
\begin{equation}\label{eq:directional-prox}
 \prox_{g_C}^{H}(v)=\Pi_C(v)+\prox_{\Psi}^{H}(v-\Pi_C(v)),\qquad
 g_C(z)=\inf_{s\in C}\Psi(z-s).
\end{equation}
The proof clips each candidate to the coordinate segments between $\Pi_C(v)$
and $v$. Both diagonal quadratic distance and the contraction-monotone penalty
cannot increase. On those segments the problem is exactly the residual
proximal problem. A dense positive metric does not support that coordinatewise
argument; the theorem is not silently extended to arbitrary metric coupling.

For one directional group let $r=v-\Pi_C(v)$, choosing $c_i$ from the active
sign of each residual. A fixed candidate peak $t\geq0$ gives
\begin{equation}\label{eq:directional-clipping}
 w_i(t)=\operatorname{sign}(r_i)\min\{|r_i|,t/c_i\},\qquad
 F(t)=at+\tfrac12bt^2+\tfrac12\sum_i h_i(|r_i|-t/c_i)_+^2.
\end{equation}
Sort the breakpoints $c_i|r_i|$ in descending order and examine a consistent active prefix $S$.
Within that interval, a stationary threshold satisfies
\begin{equation}\label{eq:directional-threshold}
 t=\frac{\sum_{i\in S}h_i|r_i|/c_i-a}
          {b+\sum_{i\in S}h_i/c_i^2}.
\end{equation}
The zero endpoint and empty active set are handled explicitly. If both prices
vanish, the residual is returned unchanged; no zero denominator is evaluated.
Sorting costs $O(k\log k)$ for a group of $k$ rows, followed by a prefix scan.
This extra proximal work can outweigh savings from a smaller principal system.

Algorithm~\ref{alg:directional} states the native operation. The local code
refuses nonrepresentable arithmetic and uses local floating-point spacing at
threshold boundaries. That treatment is important for row-scale covariance:
a fixed absolute floor in departure units would change its meaning under a
reciprocal scale transformation. Group membership remains disjoint in the
implemented interface. Overlapping groups would require a different proximal
calculation and are rejected rather than approximated as independent groups.

\begin{algorithm}[htbp]
\caption{Compute an exact directional group proximal operation in a fixed positive diagonal metric, with explicit free-price, threshold-endpoint and representability handling.}\label{alg:directional}
\begin{algorithmic}[1]
\Require Nonempty admitted interval box, disjoint groups and positive directional scales
\State Compute $p=\Pi_C(v)$ and $r=v-p$; retain hard rows separately
\For{each soft group}
\If{$a=b=0$}\State Keep the residual unchanged
\Else
\State Sort the finite breakpoints $c_i|r_i|$ in descending order
\State Find the consistent threshold interval using \eqref{eq:directional-threshold} and $t=0$
\State Set the residual coordinates using \eqref{eq:directional-clipping}
\EndIf
\EndFor
\State Return $p+w$ and continue fixed-metric relaxed splitting
\end{algorithmic}
\end{algorithm}

The joint dual condition differs from independent coordinate price caps. The
directional dual quantity is a sum of positive and negative multiplier parts
divided by their corresponding scales. Its conjugate yields one group budget.
The evaluator therefore checks every group's normalized Fenchel relation and
every hard row's complementarity separately, alongside stationarity and bounded
dual repair. A signed aggregate gap is retained diagnostically but cannot
allow an erroneous soft subgradient to cancel a negative hard-row gap.

The fixed-metric relaxation scope is $0<\alpha<2$ with the stated convexity
and attained-saddle-point assumptions. No variable-metric or arbitrary adaptive-
relaxation theorem is claimed. The complete repeated-forecast comparison uses
public SDK 0.2 and includes matched row-split, compact epigraph and direct/proximal
DAQP policies. Its primary score is \PrimaryEstimate\%, including all failed
pairs under the prespecified rule. That result refutes the registered timing
prediction on the tested population. It cannot be converted to a uniform
slowdown multiplier, nor rescued by omitting oracle limits, unsuccessful native
operations or semantically incompatible reference cells.
The physical scale controls also separate two superficially similar edits.
Multiplying an original row and its bounds by a positive conversion factor,
while dividing both directional scales by that factor, preserves the peak
and physical optimizer. Multiplying only the upper directional scale changes
the relative policy price. The teaching example intentionally exercises both
operations and compares actual native, row-split and epigraph answers. Its
hard coupling equality remains exact while the selected physical allocation
can change under the policy edit. This demonstrates the model's semantics;
it is not an additional performance experiment. The same distinction is
enforced in the research batch's unit arrays and source digests so that a
recorded scale conversion cannot silently become a new optimization objective.
Keeping those arrays attached to the original rows is essential when the
reference expands their representation into signed one-sided constraints.
The expanded reference retains a recoverable mapping to each original lower
or upper side, so its dual comparison tests the same physical quantity.
